\section*{الفصل \ref{ch:b1:quantum-introduction} --- مدخل إلى الفيزياء الكمومية}

\begin{solution}{exo:b1:quantum-introduction:1}
$h\nu$: عند \qty{100}{MHz}: $6.6 \times 10^{-26}$~جول $= \qty{4.1e-7}{eV}$؛ وعند \qty{500}{nm}:
$1240/500 = \qty{2.5}{eV}$؛ وعند \qty{0.10}{nm}: \qty{12.4}{keV}؛ وعند \qty{1}{MeV}: $\lambda =
1240/10^6 = \qty{1.2e-3}{nm} = \qty{1.2}{pm}$. وفي الليزر: $1240/633 = \qty{1.96}{eV} =
\qty{3.1e-19}{J}$، ومنه $10^{-3}/3.1 \times 10^{-19} = 3.2 \times 10^{15}$ فوتون في
الثانية.
\end{solution}

\begin{solution}{exo:b1:quantum-introduction:2}
$\lambda_0 = 1240/2.3 = \qty{540}{nm}$. وعند \qty{400}{nm} يكون $h\nu = \qty{3.1}{eV}$:
ومنه $E_{k,\max} = \qty{0.8}{eV}$ و $V_s = \qty{0.8}{V}$. ومضاعفة الشدة
تضاعف التيار وتترك $V_s$ على حاله. وعند \qty{600}{nm}
(\qty{2.07}{eV} $< W$): لا شيء، مهما تكن الشدة.
\end{solution}

\begin{solution}{exo:b1:quantum-introduction:3}
الإلكترون: $1.23/\sqrt{100} = \qty{0.12}{nm}$؛ وعند \qty{1}{keV}: \qty{0.039}{nm}.
والنترون: $E = \qty{4.0e-21}{J}$ و $p = \sqrt{2m_nE} = \qty{3.7e-24}{kg.m/s}$
و $\lambda = h/p = \qty{0.18}{nm}$. والكرة: $6.6 \times 10^{-34}/10^{-3} = \qty{7e-31}{m}$.
فللثلاثة الأولى $\lambda$ من رتبة التباعد الذري وتنعرج
(في بلورية الإلكترونات والنترونات والأشعة السينية)؛ أما الكرة فلا تنعرج أبدًا.
\end{solution}

\begin{solution}{exo:b1:quantum-introduction:4}
$E_1 = h^2/8m_eL^2 = 4.39 \times 10^{-67}/(8 \times 9.11 \times 10^{-31} \times 10^{-18}) =
\qty{6.0e-20}{J} = \qty{0.38}{eV}$ و $E_2 = 4E_1 = \qty{1.5}{eV}$؛ والفوتون
\qty{1.13}{eV}، ومنه $\lambda = 1240/1.13 = \qty{1.1}{\micro m}$ (وهو تحت أحمر قريب). والبروتون
في \qty{5}{fm}: $4.39 \times 10^{-67}/(8 \times 1.67 \times 10^{-27} \times 25 \times 10^{-30}) =
\qty{1.3e-12}{J} = \qty{8}{MeV}$ --- أي مقياس الطاقات النووية، من
الحصر وحده.
\end{solution}

\begin{solution}{exo:b1:quantum-introduction:5}
الميل $(2.05 - 0.40)/(4.0 \times 10^{14}) = \qty{4.1e-15}{V.s}$، ومنه $h = e \times 4.1 \times
10^{-15} = \qty{6.6e-34}{J.s}$. و $W = h\nu - eV_s = 2.48 - 0.40 = \qty{2.1}{eV}$،
والعتبة $\lambda_0 = 1240/2.1 = \qty{600}{nm}$: فهو السيزيوم.
\end{solution}

\begin{solution}{exo:b1:quantum-introduction:6}
$\Delta\lambda = \qty{2.43}{pm}$: ومنه $\lambda' = \qty{12.4}{pm}$. وقبل التبعثر: $1240/0.010 =
\qty{124}{keV}$؛ وبعده: $1240/0.0124 = \qty{100}{keV}$؛ فيأخذ الإلكترون
\qty{24}{keV}. وفي الضوء يكون $\Delta\lambda/\lambda = 2.4 \times 10^{-12}/5 \times 10^{-7} = 5 \times
10^{-6}$: وهو غير قابل للقياس سنة 1923 --- أما الأشعة السينية فجعلت الإزاحة أثرًا بنسبة $25\%$.
\end{solution}

\begin{solution}{exo:b1:quantum-introduction:7}
$p = \sqrt{2 \times 9.11 \times 10^{-31} \times 5 \times 10^4 \times 1.6 \times 10^{-19}} =
\qty{1.2e-22}{kg.m/s}$ و $\lambda = \qty{5.5}{pm}$ (وتصحّحها النسبية إلى
\qty{5.4}{pm})؛ وتباعد الهُدُب $\lambda D/d = 5.5 \times 10^{-12} \times 0.35/2 \times 10^{-6}
= \qty{1}{\micro m}$: وهو غير مرئي بالعين، ومن هنا \omterm{def:b1:mirrors-thin-lenses:lens}{العدسة} المكبّرة. وعشرة
إلكترونات: عشر نقاط تبدو عشوائية؛ و $10^5$: الهُدُب. فكل
إلكترون يصل عند نقطة واحدة؛ و\emph{احتمال} الوصول يتبع
موجة الشقّين --- فموجة الإلكترون الواحد تمرّ عبر الشقّين معًا.
\end{solution}

\begin{solution}{exo:b1:quantum-introduction:8}
في الذرة: $\Delta p \geq \hbar/2\Delta x = \qty{5.3e-25}{kg.m/s}$ و $E \sim \Delta p^2/2m_e =
\qty{1.5e-19}{J} \approx \qty{1}{eV}$: أي المقياس الصحيح. وفي النواة: $\Delta p \geq
\qty{1.1e-20}{kg.m/s}$؛ وكانت $\Delta p^2/2m_e$ لتكون \qty{400}{MeV}، وهي تفوق
$m_ec^2$ بكثير، فنستعمل $E \approx \Delta p\,c = \qty{3e-12}{J} = \qty{20}{MeV}$: فلا قوة
في النواة تستطيع إمساك إلكترون بهذه الطاقة --- فالنوى لا تحوي
إلكترونات (وإلكترونات بيتا تُنشَأ عند الإصدار). وفي حبة الغبار: $\Delta v \geq
\hbar/2m\Delta x = 1.05 \times 10^{-34}/(2 \times 10^{-9} \times 10^{-6}) = \qty{5e-20}{m/s}$:
وهو بلا أهمية.
\end{solution}

\begin{solution}{exo:b1:quantum-introduction:9}
(أ) $m_ev^2/r = e^2/4\pi\varepsilon_0r^2$: ومنه $v = \sqrt{e^2/4\pi\varepsilon_0m_er}$ و $E =
\tfrac12m_ev^2 - e^2/4\pi\varepsilon_0r = -e^2/8\pi\varepsilon_0r$. (ب) ولدينا $m_evr = n\hbar$ و
$m_ev^2r = e^2/4\pi\varepsilon_0$: فاقسم مربع الأولى على الثانية:
$m_er = n^2\hbar^2 \cdot 4\pi\varepsilon_0/e^2$، ومنه $r_n = n^2a_0$ و $E_n = -e^2/8\pi\varepsilon_0n^2a_0
= E_1/n^2$ مع $a_0 = \qty{52.9}{pm}$ و $E_1 = \qty{-13.6}{eV}$. (ج) و $v_1 =
\hbar/m_ea_0 = \qty{2.19e6}{m/s}$ و $v_1/c = 1/137$. (د) وفي $2 \to 1$: $\tfrac34 \times
13.6 = \qty{10.2}{eV}$، أي \qty{122}{nm}؛ وفي $3 \to 2$: $(\tfrac14 - \tfrac19) \times 13.6 =
\qty{1.89}{eV}$، أي \qty{656}{nm}. (ه) و $\lambda_n = h/m_ev_n = nh/m_ev_1 = 2\pi na_0$
(لأن $m_ev_1a_0 = \hbar$)، ومنه $2\pi r_n = 2\pi n^2a_0 = n\lambda_n$.
\end{solution}

\begin{solution}{exo:b1:quantum-introduction:10}
(أ) ليحلّ $e^2$ محلَّه المقدارُ $Ze^2$: فيصير $a_0 \to a_0/Z$ و $E_1 \to Z^2E_1$. وفي He$^+$:
\qty{54.4}{eV}. وفي U$^{91+}$: $92^2 \times 13.6 = \qty{115}{keV}$ --- لكن $v_1 = Z\alpha c
= 0.67c$: فالنموذج غير النسبي لا يعطي هناك إلا دلالة. (ب) و $a_0/207
= \qty{256}{fm}$ و $E_1 = 207 \times 13.6 = \qty{2.8}{keV}$؛ وليمان $\alpha$:
\qty{2.1}{keV} و $\lambda = \qty{0.59}{nm}$ (وهي أشعة سينية). ومدار الميون لا يزيد على
$300$ نصف قطر بروتوني: فتنزاح مستوياته انزياحًا قابلًا للقياس مع حجم البروتون،
وهكذا أُعيد قياس نصف قطر البروتون (2010). (ج) والجسيمان
يدوران حول مركزهما المشترك: \omterm{prop:b1:systems-of-points:twobody}{فالكتلة المختزلة} $m_e/2$ تنصّف
كل طاقة: ومنه $E_1 = \qty{-6.8}{eV}$، وليمان $\alpha$ عند \qty{5.1}{eV}، أي \qty{243}{nm}.
\end{solution}

\begin{solution}{exo:b1:quantum-introduction:11}
(أ) $h^2/8m_e = \qty{6.02e-38}{J.m^2}$. وعند $R = \qty{2}{nm}$: $E_e = 6.02 \times 10^{-38}/
(0.13 \times 4 \times 10^{-18}) = \qty{0.72}{eV}$ و $E_h = \qty{0.21}{eV}$، والمجموع $1.74 +
0.93 = \qty{2.67}{eV}$، أي \qty{464}{nm}: أزرق. وعند \qty{3}{nm}: الحصر بالعامل $\times4/9$،
أي \qty{2.15}{eV} و\qty{577}{nm}: أصفر. وعند \qty{4}{nm}: $\times1/4$، أي \qty{1.97}{eV} و
\qty{629}{nm}: أحمر. (ب) و\qty{520}{nm} هي \qty{2.38}{eV}: فالحصر \qty{0.65}{eV}
$= 0.93 \times (2/R)^2$، ومنه $R = \qty{2.4}{nm}$. (ج) وبلا حصر: $E_g$ وحدها،
أي \qty{713}{nm}. (د) و $\delta E = -2 \times 0.1 \times 0.41 = \qty{-0.08}{eV}$، أي $4\%$ من
\qty{2.15}{eV}: فنحو \qty{20}{nm} من التشتت --- والألوان النقية تحتاج إلى
دفعات متساوية التوزّع.
\end{solution}

\begin{solution}{exo:b1:quantum-introduction:12}
(أ) $E(\Delta x) = \hbar^2/8m\Delta x^2 + \tfrac12m\omega^2\Delta x^2$؛ وتنعدم المشتقة عند
$\Delta x^2 = \hbar/2m\omega$، حيث يتساوى الحدّان بالمقدار $\hbar\omega/4$: ومنه $E_{\min} =
\hbar\omega/2$. (ب) و $\tfrac12 \times 1.055 \times 10^{-34} \times 8.3 \times 10^{14} = \qty{4.4e-20}{J}
= \qty{0.27}{eV}$؛ و $T \sim \hbar\omega/k_B = \qty{6300}{K}$ --- فعند \omterm{def:b1:kinetic-theory:temperature}{درجة حرارة} الغرفة
يكون الاهتزاز مجمَّدًا في حالته الأساسية. (ج) و $\Delta p = \hbar/2\Delta x =
\qty{1.05e-24}{kg.m/s}$ و $E = \Delta p^2/2m_{\mathrm{He}} = 1.1 \times 10^{-48}/1.33 \times
10^{-26} = \qty{8e-23}{J} = \qty{0.5}{meV}$: أي نصف طاقة الارتباط --- فلا تستطيع الذرات
أن تسكن كفايةً لتتبلور؛ ويبقى الهيليوم سائلًا حتى الصفر
المطلق ما لم يُضغط (\qty{25}{bar}). (د) و $E_1 = \pi^2\hbar^2/2mL^2$ مقابل
$\hbar^2/2mL^2$: أي أكبر عشر مرات، كما يجب --- فالحدّ حدّ أدنى،
وحالة العلبة لها $\Delta x \approx 0.18L$ لا $L/2$.
\end{solution}

\begin{solution}{pb:b1:quantum-introduction:1}
\textbf{1.} يحرّر الضوء إلكترونات من المهبط بطاقات
مختلفة؛ وإذ يُبقى المصعد عند $-V_s$ فإنه يدفعها ولا يصل إلا ما كانت
$E_k > eV_s$. وينعدم التيار حين $eV_s = E_{k,\max}$.

\textbf{2.} $eV_s = h\nu - W$: فالميل $h/e$ والتقاطع $-W/e$؛ وتكون $V_s = 0$ عند
العتبة $\nu_0 = W/h$.

\textbf{3.} بالنقطتين الطرفيتين: $(2.36 - 0.30)/(5.0 \times 10^{14}) = \qty{4.12e-15}{V.s}$،
ومنه $h = 1.602 \times 10^{-19} \times 4.12 \times 10^{-15} = \qty{6.60e-34}{J.s}$ (وتعطي موافقة على
الخمس كلها $6.62$): أي بفارق $0.5\%$ عن $6.626 \times 10^{-34}$.

\textbf{4.} $W = h\nu - eV_s = 2.47 - 0.30 = \qty{2.2}{eV}$؛ و $\nu_0 = W/h =
\qty{5.3e14}{Hz}$ و $\lambda_0 = 1240/2.2 = \qty{560}{nm}$. أما مؤشر أحمر
(\qty{650}{nm} و\qty{1.9}{eV}) فلا ينتزع شيئًا.

\textbf{5.} تبقى $V_s$ على حالها ويتضاعف التيار: أي ضعف الفوتونات، ولكلٍّ
الطاقة نفسها. وكانت موجة لتعطي كل إلكترون طاقة أكبر مع
شدة أكبر، ولتعطي عتبة في الشدة لا في التواتر.

\textbf{6.} $10^{-6} \times 10^{-20} = \qty{1e-26}{W}$ على الذرة؛ و\qty{2.2}{eV} $=
\qty{3.5e-19}{J}$ تستغرق $3.5 \times 10^7$~ثانية --- أي سنة. ويُرصَد الإصدار
في غضون نانوثوانٍ: فالطاقة تصل كمّات كاملة.

\textbf{7.} $m_ev^2/r = e^2/4\pi\varepsilon_0r^2$: ومنه $v = \sqrt{e^2/4\pi\varepsilon_0m_er}$؛ و $E =
\tfrac12m_ev^2 - e^2/4\pi\varepsilon_0r = -e^2/8\pi\varepsilon_0r$.

\textbf{8.} $(m_evr)^2 = n^2\hbar^2$ و $m_ev^2r = e^2/4\pi\varepsilon_0$: ومنه $r_n = n^2 \cdot
4\pi\varepsilon_0\hbar^2/m_ee^2 = n^2a_0$ و $a_0 = (1.055 \times 10^{-34})^2/(8.99 \times 10^9
\times 9.11 \times 10^{-31} \times 2.57 \times 10^{-38}) = \qty{52.9}{pm}$.

\textbf{9.} $E_n = -e^2/8\pi\varepsilon_0r_n = -m_ee^4/2(4\pi\varepsilon_0)^2\hbar^2n^2$؛ ومنه $E_1 =
-2.31 \times 10^{-28}/(2 \times 5.29 \times 10^{-11}) = \qty{-2.18e-18}{J} = \qty{-13.6}{eV}$.

\textbf{10.} $v_1 = \hbar/m_ea_0 = \qty{2.19e6}{m/s}$؛ و $v_1/c = 7.3 \times 10^{-3} =
1/137$.

\textbf{11.} في $2 \to 1$: \qty{10.2}{eV} و\qty{122}{nm} (فوق بنفسجي)؛ وفي $3 \to 2$:
\qty{1.89}{eV} و\qty{656}{nm} (أحمر)؛ وفي $4 \to 2$: \qty{2.55}{eV} و\qty{486}{nm}
(أزرق مخضرّ): فخطوط بالمر هي المرئية.

\textbf{12.} $hc/\lambda = |E_1|(1/n^2 - 1/p^2)$: ومنه $R_H = |E_1|/hc = 13.6/1240 =
\qty{1.097e-2}{nm^{-1}} = \qty{1.097e7}{m^{-1}}$. وحدّ بالمر ($p \to \infty$):
$\lambda = 4/R_H = \qty{365}{nm}$.

\textbf{13.} \qty{13.6}{eV}؛ و $\lambda = 1240/13.6 = \qty{91}{nm}$.

\textbf{14.} $v_n = v_1/n$ و $\lambda_n = h/m_ev_n = nh/m_ev_1 = 2\pi na_0$ (لأن $m_ev_1a_0
= \hbar$)؛ ومنه $2\pi r_n = 2\pi n^2a_0 = n\lambda_n$: فالمدار يحمل $n$ موجة
كاملة --- وقاعدة بور قاعدة \omterm{thm:b1:wave-propagation:standing}{موجة مستقرة}.

\textbf{15.} ما يخطئ فيه: أن للإلكترون مدارًا (وكان ليشعّ؛ و
\omterm{prop:b1:quantum-introduction:box}{الحالة الأساسية} عزمها الحركي معدوم)، وأن النموذج يخفق في أيّ ذرة
فيها إلكترونان. وما يصيبه: مستويات الهيدروجين، ومن ثمّ كل خط في
طيفه. ويحلّ محلّ المدار \omterm{def:b1:quantum-introduction:wavefunction}{دالة موجية}، أي سحابة احتمال
قياسها $\sim n^2a_0$.

\textbf{16.} \omterm{thm:b1:wave-propagation:standing}{الموجات المستقرة}: $L = n\lambda/2$ و $p_n = h/\lambda_n = nh/2L$ و $E_n =
p_n^2/2m = n^2h^2/8mL^2$. وصيغة الكرة هي هي مع $L \to R$:
\omterm{prop:b1:quantum-introduction:box}{فالحالة الأساسية} تسع نصف طول موجة في نصف القطر.

\textbf{17.} $E_e = 6.02 \times 10^{-38}/(0.13 \times 4 \times 10^{-18}) = \qty{0.72}{eV}$؛
و $E_h = 0.13/0.45 \times 0.72 = \qty{0.21}{eV}$.

\textbf{18.} عند $R = \qty{2}{nm}$: $1.74 + 0.93 = \qty{2.67}{eV}$ و\qty{464}{nm}، أزرق؛
وعند \qty{3}{nm}: \qty{2.15}{eV} و\qty{577}{nm}، أصفر؛ وعند \qty{4}{nm}: \qty{1.97}{eV} و
\qty{629}{nm}، أحمر.

\textbf{19.} تنمو طاقة الحصر بنسبة $1/R^2$: فالعلبة الأصغر ترفع
المستويات، فيكون الفوتون أشدّ طاقة والضوء أزرق أكثر.

\textbf{20.} $\Delta p \sim \hbar/R = \qty{5.3e-26}{kg.m/s}$ و $E \sim \Delta p^2/2m_e^\ast =
\qty{1.2e-20}{J} = \qty{0.07}{eV}$: أي الرتبة الصحيحة (ويفسّر المعامل المضبوط
$\pi^2/2$ و $\Delta p$ الأخشن عاملَ العشرة).

\textbf{21.} $1240/577 = \qty{2.15}{eV} = \qty{3.4e-19}{J}$: ومنه $10^{-9}/3.4 \times 10^{-19}
= 2.9 \times 10^9$ فوتون في الثانية.

\textbf{22.} $10^{-3}/(1.96 \times 1.6 \times 10^{-19}) = 3.2 \times 10^{15}$ في الثانية.

\textbf{23.} $100 \times 2.48 \times 1.6 \times 10^{-19} = \qty{4e-17}{J}$ في \qty{0.1}{s}:
أي \qty{4e-16}{W}.

\textbf{24.} $h\nu = \qty{6.6e-26}{J}$: أي $1.5 \times 10^{13}$ فوتون في الثانية ---
وكلٌّ منها أصغر $10^5$ مرة من الطاقة الحرارية $k_BT$ في
الهوائي: فلا أمل، ولا حاجة؛ وصورة الموجة مضبوطة هناك.

\textbf{25.} حجم الذرات، $a_0 = \qty{53}{pm}$؛ ومقياس طاقتها،
\qty{13.6}{eV} (ومن ثمّ فوتونات \omterm{def:b1:charged-particles:eV}{بالإلكترون-فولط}، والكيمياء، والطيف المرئي)؛ و
سرعة إلكتروناتها، $\alpha c \approx \qty{2e6}{m/s}$ --- وكلها الثلاثة
من $h$ و $e$ و $m_e$ (ومن $\varepsilon_0$)، بلا ضبط أيّ شيء.
\end{solution}
